\documentclass[11pt,a4paper]{article}
\usepackage[margin=2.4cm]{geometry}
\usepackage{amsmath,booktabs,siunitx,xcolor,tabularx}
\usepackage[colorlinks=true,urlcolor=blue!60!black,citecolor=black]{hyperref}
\setlength{\parskip}{0.45em}\setlength{\parindent}{0pt}
\title{Is a parabolic flow profile a defensible idealisation for the vessels in this study?}
\author{Working note for the OCTA vessel-diameter manuscript}
\date{17 September 2026}
\begin{document}\maketitle

\section*{Why this note exists}
Section 4.7 proposes a forward model that starts from a vessel of circular cross-section carrying a
parabolic flow profile. This note checks that assumption against the vessels actually measured here,
so that the claim rests on arithmetic rather than on a rule of thumb. Everything needed to reproduce
the numbers is below.

\section*{Inputs}
\begin{tabularx}{\textwidth}{@{}l l X@{}}
\toprule
Quantity & Value used & Where it comes from \\ \midrule
Vessel diameter $d$ & \SIrange{26}{62}{\micro\metre} & ground truth, $n=7$, this study (flicker set: \SIrange{28}{48}{\micro\metre}, $n=5$) \\
Blood density $\rho$ & \SI{1060}{\kilo\gram\per\cubic\metre} & standard value for whole blood, uncited \\
Blood viscosity $\mu$ & \SI{3.5}{\milli\pascal\second}, swept \SIrange{2.5}{4.0}{\milli\pascal\second} & bulk value; see the note on apparent viscosity \cite{pries1992} \\
Mean velocity $v$ & derived, see below & total arterial flow \SI{2.82}{\micro\litre\per\minute} over six vessels \cite{zhi2012} \\
Heart rate $f$ & \SIrange{5}{10}{\hertz} (300 to 600 bpm) & swept; anaesthesia lowers heart rate and cardiac output \cite{janssen2004} \\ \bottomrule
\end{tabularx}

\paragraph{Velocity.} No measured velocity exists for these vessels, so it is derived. Zhi \textit{et al.}
measured a total arterial flow of \SI{2.82}{\micro\litre\per\minute} in the mouse retina across the six
main vessels \cite{zhi2012}. Per vessel that is $Q = 2.82/6/60 = \SI{7.83e-3}{\cubic\milli\metre\per\second}$,
and the mean velocity follows from continuity,
\begin{equation}
v = \frac{Q}{\pi (d/2)^2}.
\end{equation}
This is an estimate from a different mouse strain under a different anaesthetic, and their wording leaves it
open whether those six are six arteries or six vessels in total. If the per-vessel flow were twice what is
assumed here, every Reynolds number below would double, reaching at most $0.23$. The conclusion is
therefore also stress-tested at a velocity far above anything plausible here.

\paragraph{Apparent viscosity.} In tubes of this calibre the apparent viscosity of blood is lower than the
bulk value, the F\aa hr\ae us--Lindqvist effect, quantified over diameters from \SI{3.3}{\micro\metre} to
\SI{1978}{\micro\metre} by Pries \textit{et al.} \cite{pries1992}. A lower viscosity raises the Reynolds
number, so the sweep down to \SI{2.5}{\milli\pascal\second} is the conservative direction.

\section*{The three numbers}
\begin{align}
\text{Reynolds number} \qquad \mathrm{Re} &= \frac{\rho v d}{\mu} \\
\text{Entrance length} \qquad L_e &= \left[0.619^{1.6} + (0.0567\,\mathrm{Re})^{1.6}\right]^{1/1.6} d \\
\text{Womersley number} \qquad \alpha &= \frac{d}{2}\sqrt{\frac{2\pi f \rho}{\mu}}
\end{align}
$\mathrm{Re}$ decides laminar against turbulent. $L_e$ decides whether the parabola has room to form before
the measurement site. $\alpha$ decides whether pulsatility blunts the profile: for $\alpha \ll 1$ the flow is
quasi-steady and stays parabolic through the cardiac cycle.

\paragraph{A caution on $L_e$.} The textbook correlation $L_e \approx 0.06\,\mathrm{Re}\,d$ is not usable at
these Reynolds numbers. It is fitted at moderate $\mathrm{Re}$ and tends to zero as $\mathrm{Re}\to 0$, which
is unphysical: in creeping flow the development length tends to a constant of roughly $0.6\,d$. The
expression above carries that limit, and it is what the table uses; it is Eq.~(16) of Durst \textit{et al.},
whose constants are the $\mathrm{Re}\to 0$ and $\mathrm{Re}\to\infty$ limits of $L/D$ \cite{durst2005}.

\section*{Results}
\begin{tabular}{@{}rrrrrrr@{}}
\toprule
$d$ (\si{\micro\metre}) & $v$ (\si{\milli\metre\per\second}) & $\mathrm{Re}$ & $\mathrm{Re}$, $\mu$ \SIrange{2.5}{4.0}{\milli\pascal\second} & $L_e$ (\si{\micro\metre}) & $\alpha$ (\SI{5}{\hertz}) & $\alpha$ (\SI{10}{\hertz}) \\ \midrule
26 & 14.8 & 0.116 & 0.102--0.163 & 16 & 0.040 & 0.057 \\
30 & 11.1 & 0.101 & 0.088--0.141 & 19 & 0.046 & 0.065 \\
38 & 6.9 & 0.079 & 0.070--0.111 & 24 & 0.059 & 0.083 \\
48 & 4.3 & 0.063 & 0.055--0.088 & 30 & 0.074 & 0.105 \\
62 & 2.6 & 0.049 & 0.043--0.068 & 38 & 0.096 & 0.135 \\
\bottomrule
\end{tabular}

\paragraph{Stress test.} Forcing \SI{20}{\milli\metre\per\second} through the largest vessel, several times
the derived velocity, still gives $\mathrm{Re} = 0.38$ and an entrance length of \SI{38}{\micro\metre},
that is $0.63\,d$.

\section*{What this supports, and what it does not}
The flow in these vessels is laminar by roughly four orders of magnitude, fully developed within about two
thirds of one diameter, which is short against the spacing of bifurcations, and quasi-steady. A parabolic
profile is therefore a defensible starting point for a forward model of the OCTA signal.

Three things it does not establish. First, the profile of the \textit{OCTA signal} is not the profile of the
flow: the signal also depends on the RBC orientation artefact, on fringe washout, and on the nonlinear
relation between signal and flow speed, which is why those terms appear in the proposed model rather than flow
alone. Second, the velocity is inferred, not measured; the DyC-OCT data behind this manuscript could supply a
measured velocity for exactly these vessels and would remove that assumption. Third, a parabolic profile is
the profile of the \textit{plasma and cell mixture} treated as a continuum; it says nothing about the cell-free
layer at the wall, which is the very feature that separates the OCTA-visible red blood cell column from the
lumen boundary seen with contrast.

\begin{thebibliography}{9}
\bibitem{pries1992} A.~R. Pries, D.~Neuhaus, and P.~Gaehtgens, ``Blood viscosity in tube flow: dependence on diameter and hematocrit,'' \textit{Am. J. Physiol.} \textbf{263}(6 Pt 2), H1770--H1778 (1992). \href{https://doi.org/10.1152/ajpheart.1992.263.6.H1770}{10.1152/ajpheart.1992.263.6.H1770}
\bibitem{zhi2012} Z.~Zhi, X.~Yin, S.~Dziennis, T.~Wietecha, K.~L. Hudkins, C.~E. Alpers, and R.~K. Wang, ``Optical microangiography of retina and choroid and measurement of total retinal blood flow in mice,'' \textit{Biomed. Opt. Express} \textbf{3}(11), 2976--2986 (2012). \href{https://doi.org/10.1364/BOE.3.002976}{10.1364/BOE.3.002976}
\bibitem{durst2005} F.~Durst, S.~Ray, B.~\"Unsal, and O.~A. Bayoumi, ``The development lengths of laminar pipe and channel flows,'' \textit{J. Fluids Eng.} \textbf{127}(6), 1154--1160 (2005). Source of the $L_e$ expression, their Eq.~(16) for pipe flow.
\bibitem{janssen2004} B.~J.~A. Janssen, T.~De Celle, J.~J.~M. Debets, A.~E. Brouns, M.~F. Callahan, and T.~L. Smith, ``Effects of anesthetics on systemic hemodynamics in mice,'' \textit{Am. J. Physiol. Heart Circ. Physiol.} \textbf{287}(4), H1618--H1624 (2004). Cited for the direction of the effect only; the abstract reports cardiac output and blood pressure, not heart rate. \href{https://doi.org/10.1152/ajpheart.01192.2003}{10.1152/ajpheart.01192.2003}
\end{thebibliography}
\end{document}
